\section{Opérateurs fondamentaux du TPK : domaines, actions, invariants et
mémoire}
\label{sec:inventario-operatorio-tpk-rev11}
\index{TPK!inventaire opératoire}

\cref{def:categorias-operatorias-tpk-rev11} a d'abord fixé les huit
classes ontologico--opératoires : support et état ; sélection ; transport ;
lecteurs et quotients ; mémoire et frontière ; périodicité et retour ;
prolongation ; et sorties. Cette section définit leurs réalisations. Pour chaque
opérateur, on déclare le sextuplet \((D,C,f,I,L,M)\) du domaine, du codomaine, de la règle,
des invariants, de l'information publiée ou perdue et de la mémoire conservée. Le recensement
nominal ultérieur regroupe des symboles déjà définis par ces six données et ne
remplace ni leurs domaines ni leurs règles.

\begin{definition}[Fiche typée d'un opérateur du TPK]
Un opérateur canonique est un sextuplet
\[
 \mathcal O=(D_{\mathcal O},C_{\mathcal O},f_{\mathcal O},
 I_{\mathcal O},L_{\mathcal O},M_{\mathcal O}),
\]
où \(f_{\mathcal O}:D_{\mathcal O}\to C_{\mathcal O}\) est la règle
d'action, \(I_{\mathcal O}\) est la famille de relations qu'il conserve,
\(L_{\mathcal O}\) déclare les coordonnées publiées et la relation
d'équivalence induite par l'oubli, et \(M_{\mathcal O}\) énumère les données
qui restent dans l'état enrichi après la publication. Deux
réalisations représentent le même opérateur uniquement lorsqu'il existe une
équivalence typée qui entrelace leurs règles et leurs invariants ; partager une
période, un cardinal ou une sortie visible ne suffit pas.
\end{definition}

\subsection{Lecteurs modulaires construits}

La réduction nonadique est fixée par le représentant publié
\begin{equation}
 \rho_9(n)=1+((n-1)\bmod 9)
 =n-9\left\lfloor\frac{n-1}{9}\right\rfloor,
 \label{eq:rho9-construida-tpk-rev11}
\end{equation}
de sorte que la classe nulle s'écrit \(9\). Pour \(n\geq1\), la racine
numérique satisfait
\[
 \operatorname{dr}_9(n)=\rho_9(n),
\]
mais les deux expressions conservent des sémantiques distinctes : la première est une
section du quotient ; la seconde publie la somme itérée des chiffres. Le relèvement
tritique écrit de manière unique
\begin{equation}
 n=r_3(n)+3c_3(n),
 \qquad r_3(n)\in\{-1,0,+1\},\quad c_3(n)\in\mathbb Z,
 \label{eq:lift-mod3-tpk-rev11}
\end{equation}
et conserve simultanément le résidu orienté et le quotient de retenue.
Enfin, le lecteur de blocs décimaux est la division euclidienne
\begin{equation}
 N=1000q_{1000}(N)+r_{1000}(N),
 \qquad 0\leq r_{1000}(N)<1000,
 \label{eq:lector-mod1000-tpk-rev11}
\end{equation}
de sorte que la base \(1000\), le résidu modulo \(1000\) et la retenue \(q_{1000}\)
sont trois données liées et non un même objet.

\subsection{Classification des réalisations opératoires par domaine et fonction}

Les définitions précédentes régissent la lecture de la table. Ses quatorze
regroupements fonctionnels raffinent les huit classes opératoires et résument les
réalisations qui partagent une fonction ; une classe peut donc occuper plus
d'une ligne. La table ne sert pas de définition de remplacement. Les symboles qui
partagent un cardinal restent séparés par leur domaine et leur codomaine.

\begingroup
\footnotesize
\renewcommand{\arraystretch}{1.06}
\begin{longtable}{@{}p{.21\textwidth}p{.31\textwidth}p{.40\textwidth}@{}}
\toprule
Regroupement fonctionnel & Opérateurs et espaces & Domaine, action et invariants \\
\midrule
\endfirsthead
\toprule
Regroupement fonctionnel & Opérateurs et espaces & Domaine, action et invariants \\
\midrule
\endhead
État et mémoire de route &
\(\mathcal X_{\rm TPK}^{\rm enr}\), fibre \(\mathcal F_q\), mémoire de route &
Ils conservent position, feuillet, orientation, résidu, quotient, retenue, signature,
frontière, cylindre, préfixe, généalogie et mémoire. \\

Sélection de feuillet &
\(M_{\rm ph}:\{1,\ldots,9\}\to\{-1,0,+1\}\) &
Active l'évaluation additive, multiplicative ou le passage du seuil ; ne
déplace pas à lui seul l'état. \\

Transporte espacial &
\(T_N,T_E,T_S,T_O\), moteur réversible à deux curseurs,
\(F_{\rm obs}\) &
Réalise les translations du graphe de Cayley, met à jour le feuillet et l'orientation et
conserve chaque transition de la route. \\

Lectores modulares &
\(\rho_9\), \((r_3,c_3)\), \((q_{1000},r_{1000})\) &
Publient la phase nonadique, l'orientation tritique et le bloc décimal sans effacer les
quotients de relèvement. \\

Bloc et mémoire &
\(\mathcal K_{27}=F_{\rm obs}^{27}\), filtre de mémoire
\(\mathcal N_{27}\) &
Le premier compose vingt-sept transports ; le second sélectionne les
événements de mémoire. Un même cardinal n'implique pas un même opérateur. \\

État dodécaphasique &
\(\mathcal R_{12}\), centres \(\theta_m=30m-10\) degrés,
raffinement \(\mathcal R_{120}\) &
Organise les douze phases, leur opposition et le raffinement angulaire par incréments de
trois degrés. \\

Périodicité et profondeurs &
\(C_{12}\hookrightarrow C_{108}\twoheadrightarrow C_9\),
\(\Pi_{108}^{\rm cur}\), \(w_{108}\), \(w_{1080}\),
\(\Gamma_0(1080)\) &
Distingue la périodicité totale d'ordre \(108\), la projection des curseurs, les
mots de 108 et 1080 trits et le niveau modulaire 1080. \\

Canaux (90/120) &
masque \(C_m\), \(S_{90},S_{120}\), grammaire orbitale,
supertrace et invariants \(\nu_{120},\nu_{270}\) &
Sépare les deux généalogies, leurs queues géométriques, leurs poids et les
coordonnées qu'elles extraient de l'état. \\

Bidirectionnalité &
\(P_\pm\), inversion de route, \(\operatorname{Ext}_\pm\),
\(N_\pm\), \(\beta\), \(C_M=-\operatorname{rev}\) &
Réalise, respectivement, l'échange spectral, l'aller et le retour spatial,
la prolongation future/passée, la valence, le bilan et la conjugaison du mot. \\

Émission et coinduction &
\(L_0,L_1\), \(R_{36}\), \(G_9\), partition \(3^6=729\),
\(\varprojlim\operatorname{Cyl}_m\) &
Élève le germe, ouvre la frontière et énumère exhaustivement les
prolongations finies. La limite numérique et le relèvement enrichi
appartiennent à des types distincts. \\

Lecteurs numériques &
\(\mathscr C_\pi\), \(\mathscr P_e\), \(F_{\rm av}\) &
Construisent, par encadrements rationnels, récurrence factorielle et auto-échelle de
Fibonacci, les sections publiées de \(\pi,e,\varphi\). \\

Projection exceptionnelle &
relèvement de Witt, énumérateur de poids et ombre non sélectrice &
Transporte le même état vers l'incidence ternaire sans intervenir dans la
sélection archimédienne des chiffres. \\

Atlas et réalisation &
signature de route, atlas (81/56/13/324), involution \(C_M\), centre
électronique et opérateurs de fibre &
Transforme l'état de route en classificateurs, caractères et opérateurs avant
d'appliquer toute carte dimensionnelle. \\

Validation finie &
recensements exhaustifs, ablations et indépendance causale &
Contrôle la couverture et l'unicité à la profondeur calculée, et vérifie
qu'aucun préfixe cible n'intervient dans l'application génératrice. \\
\bottomrule
\end{longtable}
\endgroup

\Needspace{46\baselineskip}
\subsection{Index fonctionnel des opérateurs canoniques}

La classification précédente organise les réalisations par fonction. La table suivante
réunit les opérateurs après la définition catégorielle et les lecteurs
construits. Chaque ligne déclare le domaine, le codomaine, la règle et les invariants ; les
symboles permettent de reconnaître une même définition à travers ses
différentes réalisations sans identifier des opérateurs de types différents.

\begingroup
\footnotesize
\renewcommand{\arraystretch}{1.06}

\begin{longtable}{@{}r >{\raggedright\arraybackslash}p{.42\textwidth}
  >{\raggedright\arraybackslash}p{.22\textwidth}
  >{\raggedright\arraybackslash}p{.22\textwidth}@{}}
\toprule
\# & Opérateur canonique & Symbole & Capa tipada \\
\midrule
\endfirsthead
\toprule
\# & Opérateur canonique & Symbole & Capa tipada \\
\midrule
\endhead
1 & Cellule arithmétique additive--multiplicative d'ordre neuf
  & \(\mathcal A_9\) & support arithmétique \\
2 & Décomposition résiduelle et quotient nonadiques
  & \((\rho_9,q_9)\) & réduction et élévation \\
3 & Réduction ternaire orientée
  & \(\rho_3^{\rm or}\) & réduction et élévation \\
4 & Lecteur positionnel ternaire--décimal en base mille
  & \(\iota_{1000}\) & réduction et élévation \\
5 & Action cardinale orientée sur le tore arithmétique
  & \(T_d,\ d\in\{N,E,S,O\}\) & transporte cardinal \\
6 & Rétroaction locale additive--multiplicative
  & \(F_{\rm loc}\) & transporte cardinal \\
7 & Émetteur minimal à deux curseurs
  & \(T_{\min}\) & émission finie \\
8 & Coupleur décimal de signatures arithmétiques
  & \(G\) & émission finie \\
9 & Projection ternaire du catalogue
  & \(\Psi_6\) & émission finie \\
10 & Surjectivité locale du transport cardinal
  & \(\mathcal R_3\) & transporte cardinal \\
11 & Lacet de transport de vingt-sept itérations
  & \(\mathcal K_{27}\) & transporte cardinal \\
12 & État enrichi du noyau topologique de phase
  & \(\mathcal X_{\rm TPK}^{\rm enr}\) & estado enriquecido \\
13 & Transition d'état de Hensel
  & \(\mathcal H\) & estado enriquecido \\
14 & Transition exacte de cylindres ternaires--décimaux
  & \(\mathcal C_{3,1000}\) & estado enriquecido \\
15 & Signature conjointe de frontière
  & \(\Sigma_{\rm B}\) & estado enriquecido \\
16 & Relation d'extension et de survie nonadique
  & \({\rm EF}\text{--}G_9\) & prolongation et monodromie \\
17 & Monodromie nonadique avec holonomie de frontière
  & \(\mathcal M_9\) & prolongation et monodromie \\
18 & Section globale du système inverse des survivants
  & \(X_\chi=\varprojlim_m{\rm Surv}^{(\chi)}_m\)
  & prolongation et monodromie \\
19 & Catégorie bivariée de chemins arithmétiques
  & \(\mathcal P_{\rm APP}^{(2)}\) & transporte cardinal \\
20 & Chronologie observable à deux curseurs
  & \(F_{\rm obs}\) & transporte cardinal \\
21 & Alphabet décimal de quarante-deux triades
  & \(\mathcal A_{42}\) & émission finie \\
22 & Relation typée d'extension de fibres
  & \({\rm Ext}_r\) & estado enriquecido \\
23 & Transfert nonadique de fibres compatibles
  & \(K_{\rm ph}\) & prolongation et monodromie \\
24 & Raffinement des germes par avancement et rotation
  & \((P,H,Q)\) & transporte cardinal \\
25 & Cocycles de mode et d'orientation
  & \((c_{\rm mode},c_{\rm or})\) & estado enriquecido \\
26 & Prolongation triadique de vingt blocs
  & \(E_{20}\) & prolongation et monodromie \\
27 & Système dodécaphasique orienté
  & \(\mathcal R_{12}\) & lecture dodécaphasique \\
28 & Observable harmonique dodécaphasique signée
  & \(U_{12}^{\rm sgn}\) & lecture dodécaphasique \\
29 & Transformation harmonique par orbites de pas trois
  & \(\mathcal H_{12}\) & lecture dodécaphasique \\
30 & Extracteur transversal de différences cycliques
  & \(\mathcal E_{3,4}\) & lecture dodécaphasique \\
31 & Récurrence dodécaphasique de fermeture en base mille
  & \(\mathcal C_\alpha\) & fermeture en base mille \\
32 & Transducteur d'histoire enrichie en caractère angulaire
  & \((L_{\rm tick},\chi_{\rm mass})\) & interfaz dimensional \\
33 & Réseau bidirectionnel d'indice douze
  & \(\mathcal W_{\pm}\) & interfaz dimensional \\*
34 & Atlas nonadique des familles d'extension
  & \(\mathcal A_{\rm species}\) & interfaz dimensional \\
\bottomrule
\end{longtable}
\endgroup

Le symbole \(L_{\rm tick}\) désigne le lecteur d'itération et de retour. Son
codomaine est la mémoire temporelle de l'état ; il ne définit pas d'unité temporelle
physique indépendante.

La fonction sélectrice interne
\[
 M_{\rm ph}:\mathcal D_9\longrightarrow\{+1,-1,0\}
\]
appartient au TPK\@ ; son vecteur tabulé est
\[
 m_{\rm ph}:=\bigl(M_{\rm ph}(1),\ldots,M_{\rm ph}(9)\bigr)^{\mathsf T}.
\]
La fonction \(M_{\rm ph}\), le vecteur \(m_{\rm ph}\) et l'opérateur d'état
sont des objets de types distincts. La valeur scalaire \(M_{\rm ph}(g(t))\) entre
comme second argument du sélecteur \(\operatorname{Sel}_t\) ; elle ne se compose pas
directement avec l'état. Le transport cardinal \(T_d\) occupe un autre facteur
de la transition \(\mathcal U_t\) et réalise le mouvement sur APP au moyen de
son élévation \(\operatorname{Lift}(T_d)\). Cette séparation conserve un unique
conteneur, distingue la sélection de phase du transport et préserve
l'architecture complète par types.

\Needspace{9\baselineskip}
\begin{resultbox}[title={Profondeur finie et limite enrichie}]
L'énumération, les recensements et la contraction des cylindres numériques sont des
résultats exacts dans leurs domaines. L'existence et l'unicité d'une section
enrichie soumise conjointement à tous les prédicats de survie
s'obtiennent lorsque les fibres de relèvement sont non vides, unitaires et
compatibles à toute profondeur. Un calcul sur une fenêtre finie
vérifie les fibres de cette fenêtre, mais n'implique pas à lui seul le
quantificateur sur toutes les profondeurs.
\end{resultbox}

\begin{theorem}[Composition de la dynamique interne]
\label{thm:composicion-dinamica-tpk-rev11}
Une transition observable du TPK est la composition typée
\[
 \operatorname{Tra}_t(y,d):=\operatorname{Lift}(T_d)(y),
 \qquad
 \operatorname{Upd}_t:=\operatorname{Rec}_t\circ\operatorname{Mem}_t,
\]
\[
 \mathcal U_t(x)
 =\operatorname{Upd}_t\!\left(
  \operatorname{Tra}_t\!\left(
   \operatorname{Sel}_t\!\left(x,M_{\rm ph}(g(t))\right),d(t)
  \right)\right)
\]
sur \(\mathcal X_{\rm TPK}^{\rm enr}\), suivie, le cas échéant, d'un
lecteur modulaire, d'une extension cylindrique ou de l'une des deux projections.
Chaque facteur agit sur une composante distincte de l'état ; ainsi, le
sélecteur, le mouvement, l'horloge, la mémoire et la publication ne peuvent pas
se réduire à une seule réduction modulaire.
\end{theorem}

La table précédente énumère les familles opératoires. Les équivalences entre
réalisations locales doivent entrelacer leurs domaines, leurs règles et leurs invariants.
Les développements ultérieurs sélectionnent des compositions de ces opérations ;
ils n'introduisent pas un quatrième objet primitif aux côtés d'APP, TRIT et TPK.
